% Appendices follow the references, as in the institute template.
\clearpage
\phantomsection\addcontentsline{toc}{chapter}{Appendix-A}
\chapter*{Experimental Settings}\markboth{Appendix-A}{Appendix-A}
\setcounter{table}{0}\renewcommand{\thetable}{A.\arabic{table}}
\label{app:settings}

\begin{table}[ht]
\centering
\caption{Settings of the controlled study}
\label{tab:settings}
\begin{tabularx}{\textwidth}{@{}l X@{}}
\toprule
Parameter & Value \\
\midrule
UAVs $K$, base stations $B$ & 10, 3 \\
Byzantine UAVs $f$ & 3, fixed per seed across all methods \\
Rounds; synchronization period $E$ & 100; 5 \\
Local training & 1 epoch, SGD, learning rate 0.01, momentum 0.9, batch 64, weight decay $5\times10^{-4}$ \\
Model & TacticalEdgeCNN, 156{,}074 trainable parameters, no normalization layers \\
Data & Aerial-transformed CIFAR-10, Dirichlet $\alpha=0.5$, 49{,}500 training images to UAVs \\
Root set $\mathcal{D}_0$ & 500 images (50 per class), disjoint from all UAV data \\
Reputation & $\beta=0.5$, gate $\tau=0.3$, initial trust 1 \\
Staleness discount & $\exp(-0.5\,\cdot\,\text{staleness})$; late updates dropped \\
Packets & $L=1024$ parameters, 32 bits per parameter \\
Packet errors & logistic, slope $a=1$, threshold $\gamma_{\mathrm{th}}=3$\,dB \\
Round deadline; bandwidth & 2\,s; 20\,MHz \\
Handover & 3\,dB hysteresis, 2-round time to trigger \\
Trace & DeepSense~6G Scenario~23, stride 5 samples per round, segments $\ge 200$ samples \\
Reference SNR $\bar\gamma$ & assumed anchors: 15\,dB (good link), 5\,dB (weak link) \\
Jammer; parameter noise & none; none \\
Attacks & sign flip $\lambda=5$; ALIE-style $z=1.5$; IPM $\epsilon=0.5$ \\
Seeds & 41--50 \\
\bottomrule
\end{tabularx}
\end{table}

\clearpage
\phantomsection\addcontentsline{toc}{chapter}{Appendix-B}
\chapter*{Reproducing the Results}\markboth{Appendix-B}{Appendix-B}
\setcounter{table}{0}\renewcommand{\thetable}{B.\arabic{table}}
\label{app:repro}

All code, the run-level summary results and this report are in the repository
\url{https://github.com/sasank-555/uav-swarm-fl}. The raw DeepSense~6G data are
access-controlled and must be obtained from the dataset provider; per-run model
checkpoints are kept locally and are not distributed. The training-source commit
pinned for the controlled study is recorded in
\texttt{results/CONTEXT\_STUDY\_FINAL.md}.

\begin{table}[ht]
\centering
\caption{Repository layout}
\label{tab:repo}
\begin{tabularx}{\textwidth}{@{}l X@{}}
\toprule
Path & Contents \\
\midrule
\texttt{fl/} & Simulator: data, model, client, channel, topology, attacks, aggregators, server \\
\texttt{tests/} & Automated simulator, coordinator and analysis tests \\
\texttt{run\_context\_study.py} & Plans and runs the \NumRuns-run controlled study \\
\texttt{analyze\_context\_study.py} & Prespecified analysis with paired bootstrap \\
\texttt{results/context\_study/} & Committed summary CSV and analysis of the controlled study \\
\texttt{results/context\_final/} & Local output of the final-profile reproduction command \\
\texttt{results/base\_paper\_comparison/} & Compact evidence, paired estimates and validation of the 240-run follow-up \\
\texttt{scripts/make\_paper\_figures.py} & Builds every figure, table and number in the paper and this report \\
\texttt{paper/} & Manuscript in IEEE letter format \\
\texttt{report/} & This report \\
\bottomrule
\end{tabularx}
\end{table}

Run these commands from the repository root. The final profile writes to
\path{results/context_final/}. The published compact summaries are in
\path{results/context_study/}.
\begin{lstlisting}
pip install -r requirements.txt
python prepare_data.py
python scripts/prepare_deepsense23.py /path/to/Scenario23 --out data/channel_traces/deepsense23.npz
python run_context_study.py --profile final --device cuda --parallel 4 --out-dir results/context_final
python analyze_context_study.py results/context_final
\end{lstlisting}
To rebuild figures and numbers from the committed summaries, then compile the
letter and report, use:
\begin{lstlisting}
python scripts/make_paper_figures.py
(cd paper && tectonic -X compile main.tex)
(cd report && tectonic -X compile main.tex)
\end{lstlisting}
These last three commands require no new training. To plot a fresh reproduction
instead, pass \lstinline|--csv results/context_final/context_runs.csv| to the
figure script; its default is the committed
\texttt{results/context\_study/context\_runs.csv}.

The completed base-paper follow-up uses frozen training revision
\texttt{ad06055} and the same runner with
\lstinline|--aggs bp_fresh bp_reuse|. Its full checkpoints, coordinator audit and
prepared inputs remain local. Compact publication checks source/data fingerprints,
terminal histories and all 120 paired client-partition summaries. Recompute the
follow-up intervals from the committed summary data without training:
\begin{lstlisting}
python scripts/analyze_base_paper.py --csv results/base_paper_comparison/base_paper_runs.csv
\end{lstlisting}
Then repeat the figure-generation and document-build commands above. The result
CSV records the final attempt's runtime
profile and resumed round; an attempt that resumed may include earlier work from a
different device. The source fingerprints of both studies are retained separately.
